%% kapa: exam=1 round=26 year=2015 period=2 form=A subject=accounting subject_ko=회계학 seq=1 number=41
%% source: Q-Net 감정평가사 기출문제 (cst00309) artlSeq=5138001 file=extracted/1008/5138001/제26회_감정평가사 1차_기출/제26회_감정평가사 1차_기출/2교시_감평_A형.pdf
%% conversion: kordoc 4.24.51 -> tools/gen1.py (auto); review=proofread-sonnet; flags=none
\begin{question}{41}
다음은 공정가치에 관한 정의이다. 괄호 안에 들어갈 적절한 단어를 나타낸 것은?
\begin{framedtext}
공정가치는 원칙적으로 측정일에 ( ㄱ ) 사이의 정상거래에서 자산을 매도하면서 수취하거나 ( ㄴ )을(를) 이전하면서 지급하게 될 ( ㄷ )을 의미한다.
\end{framedtext}

\noindent\hspace*{2em}\makebox[8em][l]{(ㄱ)}\makebox[4em][l]{(ㄴ)}(ㄷ)\par
\begin{choices}
  \item \makebox[8em][l]{매도자와 매수자}\makebox[4em][l]{손실}유출가격
  \item \makebox[8em][l]{이해관계자}\makebox[4em][l]{손실}유입가격
  \item \makebox[8em][l]{이해관계자}\makebox[4em][l]{부채}유입가격
  \item \makebox[8em][l]{시장참여자}\makebox[4em][l]{부채}유입가격
  \item \makebox[8em][l]{시장참여자}\makebox[4em][l]{부채}유출가격
\end{choices}
\end{question}
\answer{5}
